\section{附录：文件列表与核心代码}

本文计算程序和结果表均存放于 \texttt{A题} 文件夹下。为便于复现，表~\ref{tab:appendix_files}
列出本文提交和复核所需的主要文件。其中 \texttt{problem1.py} 至 \texttt{problem4.py}
对应四个子问题的主程序，两个 analytic 文件用于线性情形的解析解校验和功率复核。

\begin{table}[H]
  \centering
  \caption{主要程序文件和结果文件}
  \label{tab:appendix_files}
  \small
  \begin{tabular}{p{0.34\textwidth}p{0.56\textwidth}}
    \toprule
    文件名 & 作用说明 \\
    \midrule
    \texttt{problem1.py} & 求解问题一中线性阻尼和非线性阻尼下的二自由度垂荡响应。\\
    \texttt{problem2.py} & 对问题二的线性阻尼系数及非线性阻尼参数进行功率优化。\\
    \texttt{problem3.py} & 求解问题三中垂荡--纵摇四自由度非线性几何耦合模型。\\
    \texttt{problem4.py} & 对问题四的垂荡阻尼和旋转阻尼进行联合优化。\\
    \texttt{problem1\_analytic\_check.py} & 构造问题一线性阻尼情形的时域解析解，并与数值解比较。\\
    \texttt{problem2\_linear\_analytic.py} & 利用复振幅法计算问题二线性阻尼的稳态平均功率和最优阻尼。\\
    \texttt{result1-1.xlsx} & 问题一线性阻尼情形的计算结果表。\\
    \texttt{result1-2.xlsx} & 问题一非线性阻尼情形的计算结果表。\\
    \texttt{result3.xlsx} & 问题三给定阻尼参数下的四自由度响应结果表。\\
    \bottomrule
  \end{tabular}
\end{table}

以下仅给出各程序中与建模和优化直接相关的核心代码片段。为控制附录篇幅，参数定义、文件读写、
绘图语句和重复性辅助代码均已省略。

\subsection{问题一：二自由度垂荡响应}

\begin{lstlisting}[language=Python,caption={问题一垂荡运动方程与数值积分核心片段}]
def ode_system(t, x, c_damp, nonlinear=False, alpha=0.5):
    z_b, zb_dot, z_z, zz_dot = x
    rel_disp = z_b - z_z
    rel_vel = zb_dot - zz_dot
    if nonlinear:
        f_damp = c_damp * abs(rel_vel) ** alpha * rel_vel
    else:
        f_damp = c_damp * rel_vel
    f_pto = k_spring * rel_disp + f_damp
    zb_ddot = (f_amp*np.cos(omega*t) - C_wave*zb_dot
               - K_hydro*z_b - f_pto) / (m_b + m_add)
    zz_ddot = f_pto / m_z
    return [zb_dot, zb_ddot, zz_dot, zz_ddot]
\end{lstlisting}

\subsection{问题二：阻尼参数优化}

\begin{lstlisting}[language=Python,caption={问题二平均功率与非线性阻尼优化核心片段}]
def damping_power(c, alpha, rv):
    return c * np.abs(rv) ** (alpha + 2)

def search_nonlinear_surface(c_range=(1e-2, 100000), alpha_range=(0.0, 1.0),
                             n_c=36, n_alpha=31):
    log_bounds = (np.log10(c_range[0]), np.log10(c_range[1]))
    log_grid = np.linspace(log_bounds[0], log_bounds[1], n_c)
    alpha_grid = np.linspace(alpha_range[0], alpha_range[1], n_alpha)
    records, best = [], (-np.inf, None, None)
    for a in alpha_grid:
        for log_c in log_grid:
            c = 10 ** log_c
            p_val = rk4_power(c, a, n_total=32, n_drop=16)
            records.append((c, a, p_val))
            if p_val > best[0]:
                best = (p_val, c, a)

    def objective(x):
        log_c, alpha = x
        c = 10 ** log_c
        return -rk4_power(c, alpha, n_total=36, n_drop=18)

    x0 = np.array([np.log10(best[1]), best[2]])
    powell = minimize(objective, x0, method="Powell",
                      bounds=[log_bounds, alpha_range])
    de = differential_evolution(objective, bounds=[log_bounds, alpha_range])
    # 对粗扫、Powell 和差分进化的候选点统一长时间复算。
\end{lstlisting}

\subsection{问题三：四自由度非线性耦合模型}

\begin{lstlisting}[language=Python,caption={问题三非线性几何耦合方程与 Radau 积分核心片段}]
def ode_system(t, x):
    z_z, z_b, th_z, th_b, zz_dot, zb_dot, thz_dot, thb_dot = x
    angle = th_z + th_b
    sin_a = np.sin(angle)
    cos_a = np.cos(angle)
    thb_ddot = (-K_pitch*th_b - C_pitch*thb_dot
                + c_rotary*thz_dot + k_torsion*th_z
                + L_amp*np.cos(omega*t)) / (I_b + J_add)
    Iz_eff = I_z + m_z * z_z**2
    thz_ddot = (-c_rotary*thz_dot - k_torsion*th_z
                + m_z*g*z_z*sin_a) / Iz_eff - thb_ddot
    m_eff = m_b + m_add + m_z * sin_a**2
    zb_ddot = (
        k_linear*(z_z-l0+m_z*g/k_linear)*cos_a + c_linear*zz_dot*cos_a
        - C_heave*zb_dot + f_amp*np.cos(omega*t) - K_heave*z_b
        + sin_a*(m_z*(z_z*thz_ddot + 2*zz_dot*thz_dot
        + z_z*thb_ddot + 2*zz_dot*thb_dot) - m_z*g*sin_a)
    ) / m_eff
    zz_ddot = (
        -k_linear*(z_z-l0) - c_linear*zz_dot - m_z*g*cos_a
        - m_z*(-z_z*thz_dot**2 + zb_ddot*cos_a
        - z_z*thb_dot**2 - 2*z_z*thb_dot*thz_dot)
    ) / m_z
    return [zz_dot, zb_dot, thz_dot, thb_dot,
            zz_ddot, zb_ddot, thz_ddot, thb_ddot]

sol = solve_ivp(ode_system, (0, 40*T_period), x0, t_eval=t_eval,
                method="Radau", rtol=1e-8, atol=1e-10)
\end{lstlisting}

\subsection{问题四：垂荡与旋转阻尼联合优化}

\begin{lstlisting}[language=Python,caption={问题四平均功率计算和局部优化核心片段}]
def average_power(c_linear, c_rotary, n_periods=45, drop_periods=25):
    sol = solve_response(c_linear, c_rotary, n_periods=n_periods)
    t = sol.t
    mask = t >= drop_periods * T_period
    zz_dot = sol.y[4, mask]
    thz_dot = sol.y[6, mask]
    t = t[mask]
    p_heave = c_linear * zz_dot ** 2
    p_pitch = c_rotary * thz_dot ** 2
    p_linear = np.trapezoid(p_heave, t) / (t[-1] - t[0])
    p_rotary = np.trapezoid(p_pitch, t) / (t[-1] - t[0])
    return PowerResult(c_linear, c_rotary, p_linear+p_rotary,
                       p_linear, p_rotary)

def coarse_scan():
    values = np.linspace(0.0, 100000.0, 11)
    rows = []
    for c1 in values:
        for c2 in values:
            rows.append(average_power(c1, c2).__dict__)
    return pd.DataFrame(rows).sort_values("p_total", ascending=False)

def refine_from(best):
    cache = {}
    def objective(v):
        c1, c2 = v
        key = (round(c1), round(c2))
        if key not in cache:
            cache[key] = average_power(c1, c2)
        return -cache[key].p_total
    opt = minimize(objective, [best.c_linear, best.c_rotary],
                   method="Powell", bounds=[(0, 1e5), (0, 1e5)])
    return average_power(opt.x[0], opt.x[1])
\end{lstlisting}

\subsection{线性问题的解析校验}

\begin{lstlisting}[language=Python,caption={问题一线性阻尼时域解析解核心片段}]
def analytic_solution(times):
    A, b = state_matrices()
    q = np.linalg.solve(1j*omega*np.eye(4) - A, b)
    x0 = np.zeros(4)
    x_ss_0 = np.real(q)
    exact = []
    for t in times:
        x_steady = np.real(q * np.exp(1j*omega*t))
        x_exact = x_steady + expm(A*t).dot(x0 - x_ss_0)
        exact.append(x_exact)
    return np.asarray(exact), q
\end{lstlisting}

\begin{lstlisting}[language=Python,caption={问题二线性阻尼稳态功率解析优化核心片段}]
def complex_amplitude(c):
    M = np.array([[m_b + m_add, 0], [0, m_z]], dtype=complex)
    C = np.array([[C_w + c, -c], [-c, c]], dtype=complex)
    K = np.array([[K_h + k_s, -k_s], [-k_s, k_s]], dtype=complex)
    F = np.array([f_amp, 0], dtype=complex)
    return np.linalg.solve(K - omega**2*M + 1j*omega*C, F)

def analytic_power(c):
    z = complex_amplitude(c)
    return 0.5 * c * omega**2 * abs(z[0] - z[1]) ** 2
\end{lstlisting}
